\section{Attack Asymmetry, Shared Signals y Zero Trust}

Al atacante le basta con encontrar una credential, una session, un help-desk o un integration gap; el equipo defensor, en cambio, debe coordinar endpoint, identity, application, data y human workflow. En clase se vinculó esta asimetría con los obstáculos para compartir información: el diseño de sistemas necesita defense in depth y también intercambiar risk signal dentro de los límites de la privacidad y la ley, para reducir el tiempo que cada organización tarda en descubrir por su cuenta el mismo patrón de ataque.

\subsection{Una sola debilidad frente a varias capas de defensa}

El equipo defensor no puede cargar todo el riesgo sobre el IdP. Endpoint compromise, phishing, malware, OAuth consent, stale account, over-privileged service account y application bug pueden, cada uno, eludir una verificación única. La ventaja de la identity layer es que ve el inicio de sesión entre aplicaciones y el estado de factor, session y group, pero aun así necesita recibir indicadores de EDR, network, application y customer-reported indicators.

\lecturefigure{10-shared-defense.png}{El atacante necesita un solo hueco; el defensor debe coordinar controles en varias capas y señales entre organizaciones.}{Subtítulos locales 21:20--24:10; redibujo conceptual.}

\begin{knowledgebox}{Lectura de la figura: la shared signal también debe minimizarse}
Un indicator puede ser una suspicious IP, session risk, credential compromise o device posture; antes de compartirlo deben definirse schema, confidence, purpose, retention y recipient. Si se transmite muy poco, se pierde valor defensivo; si se transmite demasiado, aumenta el riesgo para datos privados y comercialmente sensibles. La signal debe activar una policy evaluation local, no otorgar a la fuente externa la facultad de emitir directamente el deny final.
\end{knowledgebox}

\subsection{Zero Trust: confianza explícita, contextual y revocable}

\term{Zero Trust} no significa «nunca confiar en nadie», sino no conceder confianza permanente por la ubicación en la red o por un solo inicio de sesión. NIST SP 800-207 enfatiza que cada acceso se decida de forma explícita según subject, asset, environment y policy, que se aplique least privilege y que se suponga que la red pudo haber sido comprometida. Las identity policy modernas, además, reciben señales de riesgo de forma continua durante la session y aplican step-up, reducción de privilegios o logout.

\lecturefigure{11-zero-trust-policy.png}{Zero Trust convierte el contexto de identidad, dispositivo, riesgo y recurso en una policy continua y revocable.}{NIST SP 800-207; Okta policy/Identity Threat Protection; redibujo conceptual.}

\begin{knowledgebox}{Lectura de la figura: la policy evaluation no ocurre solo al iniciar sesión}
Identify establece el principal y el authenticator; evaluate reúne device, network y risk; enforce produce least privilege, step-up o deny; re-evaluate aplica revoke o universal logout cuando cambia el riesgo de la session. La evaluación continua no implica vigilancia ilimitada, sino elegir las señales necesarias para cada action de alto valor y dar a la session una expiry y una revocation path claras.
\end{knowledgebox}

\begin{warningbox}{Zero Trust no corrige un entitlement desordenado}
Si la organización no sabe quién es dueño de cada service account, qué permisos corresponden a cada group, cuándo se propaga una baja de personal o si una aplicación admite revoke, autenticar con más frecuencia solo sirve para verificar una y otra vez al mismo sujeto con privilegios excesivos. Zero Trust presupone inventory, ownership, least privilege y lifecycle hygiene.
\end{warningbox}

\subsection{Resumen de la sección}

La attack asymmetry exige defensa en varias capas y señales compartidas; Zero Trust, a su vez, convierte esas señales en decisiones de acceso explícitas, de mínimo privilegio y revocables. Ambas dependen de un identity state limpio y de límites de auditoría claros.
